\documentclass[11pt, letterpaper]{article}
% ============================================================
% preamble.tex — Shared LaTeX preamble for Learning to Autolens
% ============================================================
% Contains: packages, physics notation macros, formatting.
% Include in each module via % ============================================================
% preamble.tex — Shared LaTeX preamble for Learning to Autolens
% ============================================================
% Contains: packages, physics notation macros, formatting.
% Include in each module via % ============================================================
% preamble.tex — Shared LaTeX preamble for Learning to Autolens
% ============================================================
% Contains: packages, physics notation macros, formatting.
% Include in each module via \input{../preamble}
% ============================================================

% --- Packages ---
\usepackage[utf8]{inputenc}
\usepackage[T1]{fontenc}
\usepackage{amsmath, amssymb, amsthm}
\usepackage{physics}       % \vb, \grad, \div, \curl, \dd, etc.
\usepackage{mathtools}     % \coloneqq, \prescript
\usepackage{bm}            % \bm for bold math
\usepackage{graphicx}
\usepackage{float}
\usepackage{booktabs}      % Professional tables
\usepackage{hyperref}
\usepackage{cleveref}      % \cref for smart cross-refs
\usepackage{xcolor}
\usepackage[margin=1in]{geometry}
\usepackage{enumitem}      % Custom list formatting
\usepackage{fancyhdr}      % Headers and footers
\usepackage{tcolorbox}     % Colored boxes for key results
\usepackage{listings}      % Code listings

% --- Color scheme ---
\definecolor{codeblue}{RGB}{31, 119, 180}
\definecolor{codegray}{RGB}{128, 128, 128}
\definecolor{keyresult}{RGB}{39, 123, 69}
\definecolor{exercise}{RGB}{178, 102, 0}

% --- tcolorbox environments ---
\newtcolorbox{keyresult}[1][]{
  colback=keyresult!5, colframe=keyresult!70,
  fonttitle=\bfseries, title={Key Result}, #1
}

\newtcolorbox{pythonbox}[1][]{
  colback=codeblue!5, colframe=codeblue!50,
  fonttitle=\bfseries, title={PyAutoLens Implementation}, #1
}

\newtcolorbox{exercisebox}[1][]{
  colback=exercise!5, colframe=exercise!50,
  fonttitle=\bfseries, title={Exercise}, #1
}

% --- Code listing style ---
\lstset{
  language=Python,
  basicstyle=\ttfamily\small,
  keywordstyle=\color{codeblue}\bfseries,
  commentstyle=\color{codegray}\itshape,
  stringstyle=\color{keyresult},
  breaklines=true,
  frame=single,
  framerule=0.5pt,
  rulecolor=\color{codegray!50},
}

% --- Physics macros ---
% Vectors (bold upright)
\newcommand{\vtheta}{\bm{\theta}}       % Image-plane position
\newcommand{\vbeta}{\bm{\beta}}         % Source-plane position
\newcommand{\valpha}{\bm{\alpha}}       % Deflection angle
\newcommand{\vx}{\bm{x}}               % Generic position

% Lensing quantities
\newcommand{\thetaE}{\theta_{\rm E}}    % Einstein radius
\newcommand{\SigmaCr}{\Sigma_{\rm cr}}  % Critical surface density
\newcommand{\Dd}{D_{\rm d}}             % Angular diameter distance to lens
\newcommand{\Ds}{D_{\rm s}}             % Angular diameter distance to source
\newcommand{\Dds}{D_{\rm ds}}           % Angular diameter distance lens-source

% Statistical quantities
\newcommand{\chisq}{\chi^2}             % Chi-squared
\newcommand{\chisqred}{\chi^2_{\rm red}} % Reduced chi-squared
\newcommand{\Like}{\mathcal{L}}         % Likelihood
\newcommand{\Evid}{\mathcal{Z}}         % Evidence
\newcommand{\Post}{\mathcal{P}}         % Posterior

% Abbreviations for references
\newcommand{\CK}{C\&K}                  % Congdon & Keeton (2018)
\newcommand{\NB}{N\&B}                  % Narayan & Bartelmann (1997)
\newcommand{\Night}{Nightingale+18}     % Nightingale, Dye & Massey (2018)

% --- Headers ---
\pagestyle{fancy}
\fancyhf{}
\fancyhead[L]{\small\textit{Learning to Autolens}}
\fancyhead[R]{\small\thepage}
\renewcommand{\headrulewidth}{0.4pt}

% --- Theorem environments ---
\theoremstyle{definition}
\newtheorem{definition}{Definition}[section]
\newtheorem{example}{Example}[section]

% ============================================================

% --- Packages ---
\usepackage[utf8]{inputenc}
\usepackage[T1]{fontenc}
\usepackage{amsmath, amssymb, amsthm}
\usepackage{physics}       % \vb, \grad, \div, \curl, \dd, etc.
\usepackage{mathtools}     % \coloneqq, \prescript
\usepackage{bm}            % \bm for bold math
\usepackage{graphicx}
\usepackage{float}
\usepackage{booktabs}      % Professional tables
\usepackage{hyperref}
\usepackage{cleveref}      % \cref for smart cross-refs
\usepackage{xcolor}
\usepackage[margin=1in]{geometry}
\usepackage{enumitem}      % Custom list formatting
\usepackage{fancyhdr}      % Headers and footers
\usepackage{tcolorbox}     % Colored boxes for key results
\usepackage{listings}      % Code listings

% --- Color scheme ---
\definecolor{codeblue}{RGB}{31, 119, 180}
\definecolor{codegray}{RGB}{128, 128, 128}
\definecolor{keyresult}{RGB}{39, 123, 69}
\definecolor{exercise}{RGB}{178, 102, 0}

% --- tcolorbox environments ---
\newtcolorbox{keyresult}[1][]{
  colback=keyresult!5, colframe=keyresult!70,
  fonttitle=\bfseries, title={Key Result}, #1
}

\newtcolorbox{pythonbox}[1][]{
  colback=codeblue!5, colframe=codeblue!50,
  fonttitle=\bfseries, title={PyAutoLens Implementation}, #1
}

\newtcolorbox{exercisebox}[1][]{
  colback=exercise!5, colframe=exercise!50,
  fonttitle=\bfseries, title={Exercise}, #1
}

% --- Code listing style ---
\lstset{
  language=Python,
  basicstyle=\ttfamily\small,
  keywordstyle=\color{codeblue}\bfseries,
  commentstyle=\color{codegray}\itshape,
  stringstyle=\color{keyresult},
  breaklines=true,
  frame=single,
  framerule=0.5pt,
  rulecolor=\color{codegray!50},
}

% --- Physics macros ---
% Vectors (bold upright)
\newcommand{\vtheta}{\bm{\theta}}       % Image-plane position
\newcommand{\vbeta}{\bm{\beta}}         % Source-plane position
\newcommand{\valpha}{\bm{\alpha}}       % Deflection angle
\newcommand{\vx}{\bm{x}}               % Generic position

% Lensing quantities
\newcommand{\thetaE}{\theta_{\rm E}}    % Einstein radius
\newcommand{\SigmaCr}{\Sigma_{\rm cr}}  % Critical surface density
\newcommand{\Dd}{D_{\rm d}}             % Angular diameter distance to lens
\newcommand{\Ds}{D_{\rm s}}             % Angular diameter distance to source
\newcommand{\Dds}{D_{\rm ds}}           % Angular diameter distance lens-source

% Statistical quantities
\newcommand{\chisq}{\chi^2}             % Chi-squared
\newcommand{\chisqred}{\chi^2_{\rm red}} % Reduced chi-squared
\newcommand{\Like}{\mathcal{L}}         % Likelihood
\newcommand{\Evid}{\mathcal{Z}}         % Evidence
\newcommand{\Post}{\mathcal{P}}         % Posterior

% Abbreviations for references
\newcommand{\CK}{C\&K}                  % Congdon & Keeton (2018)
\newcommand{\NB}{N\&B}                  % Narayan & Bartelmann (1997)
\newcommand{\Night}{Nightingale+18}     % Nightingale, Dye & Massey (2018)

% --- Headers ---
\pagestyle{fancy}
\fancyhf{}
\fancyhead[L]{\small\textit{Learning to Autolens}}
\fancyhead[R]{\small\thepage}
\renewcommand{\headrulewidth}{0.4pt}

% --- Theorem environments ---
\theoremstyle{definition}
\newtheorem{definition}{Definition}[section]
\newtheorem{example}{Example}[section]

% ============================================================

% --- Packages ---
\usepackage[utf8]{inputenc}
\usepackage[T1]{fontenc}
\usepackage{amsmath, amssymb, amsthm}
\usepackage{physics}       % \vb, \grad, \div, \curl, \dd, etc.
\usepackage{mathtools}     % \coloneqq, \prescript
\usepackage{bm}            % \bm for bold math
\usepackage{graphicx}
\usepackage{float}
\usepackage{booktabs}      % Professional tables
\usepackage{hyperref}
\usepackage{cleveref}      % \cref for smart cross-refs
\usepackage{xcolor}
\usepackage[margin=1in]{geometry}
\usepackage{enumitem}      % Custom list formatting
\usepackage{fancyhdr}      % Headers and footers
\usepackage{tcolorbox}     % Colored boxes for key results
\usepackage{listings}      % Code listings

% --- Color scheme ---
\definecolor{codeblue}{RGB}{31, 119, 180}
\definecolor{codegray}{RGB}{128, 128, 128}
\definecolor{keyresult}{RGB}{39, 123, 69}
\definecolor{exercise}{RGB}{178, 102, 0}

% --- tcolorbox environments ---
\newtcolorbox{keyresult}[1][]{
  colback=keyresult!5, colframe=keyresult!70,
  fonttitle=\bfseries, title={Key Result}, #1
}

\newtcolorbox{pythonbox}[1][]{
  colback=codeblue!5, colframe=codeblue!50,
  fonttitle=\bfseries, title={PyAutoLens Implementation}, #1
}

\newtcolorbox{exercisebox}[1][]{
  colback=exercise!5, colframe=exercise!50,
  fonttitle=\bfseries, title={Exercise}, #1
}

% --- Code listing style ---
\lstset{
  language=Python,
  basicstyle=\ttfamily\small,
  keywordstyle=\color{codeblue}\bfseries,
  commentstyle=\color{codegray}\itshape,
  stringstyle=\color{keyresult},
  breaklines=true,
  frame=single,
  framerule=0.5pt,
  rulecolor=\color{codegray!50},
}

% --- Physics macros ---
% Vectors (bold upright)
\newcommand{\vtheta}{\bm{\theta}}       % Image-plane position
\newcommand{\vbeta}{\bm{\beta}}         % Source-plane position
\newcommand{\valpha}{\bm{\alpha}}       % Deflection angle
\newcommand{\vx}{\bm{x}}               % Generic position

% Lensing quantities
\newcommand{\thetaE}{\theta_{\rm E}}    % Einstein radius
\newcommand{\SigmaCr}{\Sigma_{\rm cr}}  % Critical surface density
\newcommand{\Dd}{D_{\rm d}}             % Angular diameter distance to lens
\newcommand{\Ds}{D_{\rm s}}             % Angular diameter distance to source
\newcommand{\Dds}{D_{\rm ds}}           % Angular diameter distance lens-source

% Statistical quantities
\newcommand{\chisq}{\chi^2}             % Chi-squared
\newcommand{\chisqred}{\chi^2_{\rm red}} % Reduced chi-squared
\newcommand{\Like}{\mathcal{L}}         % Likelihood
\newcommand{\Evid}{\mathcal{Z}}         % Evidence
\newcommand{\Post}{\mathcal{P}}         % Posterior

% Abbreviations for references
\newcommand{\CK}{C\&K}                  % Congdon & Keeton (2018)
\newcommand{\NB}{N\&B}                  % Narayan & Bartelmann (1997)
\newcommand{\Night}{Nightingale+18}     % Nightingale, Dye & Massey (2018)

% --- Headers ---
\pagestyle{fancy}
\fancyhf{}
\fancyhead[L]{\small\textit{Learning to Autolens}}
\fancyhead[R]{\small\thepage}
\renewcommand{\headrulewidth}{0.4pt}

% --- Theorem environments ---
\theoremstyle{definition}
\newtheorem{definition}{Definition}[section]
\newtheorem{example}{Example}[section]


\fancyhead[C]{\small\textbf{Module 06: Multi-Component Mass Models}}

\title{\textbf{Module 06: Multi-Component Mass Models}\\[0.5em]
\large Theory Companion --- Learning to Autolens}
\author{Rodrigo C\'ordova Rosado\\
\small Harvard-Smithsonian Center for Astrophysics\\
\small \texttt{rodrigo.cordova\_rosado@cfa.harvard.edu}}
\date{\today}

\begin{document}
\maketitle

\begin{abstract}
This document covers the theory of decomposed (stellar + dark matter) mass
models for strong gravitational lenses. We derive the mass-to-light ratio
scaling, NFW convergence, the bulge-halo conspiracy, and the dark matter
fraction. References: \CK\ Ch.~4; Treu \& Koopmans (2004);
Navarro, Frenk \& White (1996).
\end{abstract}

\tableofcontents
\newpage

% ============================================================
\section{Mass Decomposition}
\label{sec:decomposition}
% ============================================================

The total convergence of a lens galaxy is the sum of its components:
\begin{keyresult}
\begin{equation}
\kappa_{\rm total}(\vtheta) = \kappa_*(\vtheta) + \kappa_{\rm DM}(\vtheta)
\label{eq:kappa-total}
\end{equation}
\end{keyresult}

\noindent where $\kappa_*$ is the stellar convergence and $\kappa_{\rm DM}$ is
the dark matter convergence. The deflection angles add by superposition:
$\valpha_{\rm total} = \valpha_* + \valpha_{\rm DM} + \valpha_{\rm shear}$.

% ============================================================
\section{Stellar Mass: Light Traces Mass}
\label{sec:stellar}
% ============================================================

The stellar mass surface density follows the light distribution:
\begin{equation}
\Sigma_*(\vtheta) = \Upsilon_* \cdot I(\vtheta)
\label{eq:mass-to-light}
\end{equation}
where $\Upsilon_*$ is the stellar mass-to-light ratio ($M_\odot/L_\odot$).
The stellar convergence is:
\begin{equation}
\kappa_*(\vtheta) = \frac{\Upsilon_* \cdot I(\vtheta)}{\SigmaCr}
\label{eq:kappa-stellar}
\end{equation}

For a S\'ersic profile with index $n$ and effective radius $R_e$, the
enclosed mass within radius $\theta$ is:
\begin{equation}
M_*(<\theta) = 2\pi \Upsilon_* \int_0^\theta I(\theta')\,\theta'\,d\theta'
= \Upsilon_* \cdot L(<\theta)
\end{equation}

Typical values: $\Upsilon_* \sim 2$--$8\,M_\odot/L_\odot$ for ellipticals
(depends on stellar age, metallicity, and IMF).

% ============================================================
\section{Dark Matter: NFW Profile}
\label{sec:nfw}
% ============================================================

\subsection{3D Density}

The NFW profile (Navarro et al.\ 1996):
\begin{equation}
\rho(r) = \frac{\rho_s}{(r/r_s)(1+r/r_s)^2}
\label{eq:nfw-3d}
\end{equation}

The halo is characterized by the virial mass $M_{200}$ (mass within the
radius where the mean density is 200$\times$ the critical density) and
concentration $c = r_{200}/r_s$.

\subsection{Projected Convergence}

The 2D convergence for an NFW halo (Bartelmann 1996; \CK\ eq.~4.62):
\begin{equation}
\kappa_{\rm NFW}(x) = \frac{2\rho_s r_s}{\SigmaCr} \cdot f(x)
\label{eq:nfw-kappa}
\end{equation}
where $x = \theta/\theta_s$ ($\theta_s = r_s/\Dd$) and:
\begin{equation}
f(x) = \begin{cases}
\frac{1}{x^2-1}\left(1 - \frac{1}{\sqrt{1-x^2}}\text{arccosh}\frac{1}{x}\right)
& x < 1 \\[6pt]
\frac{1}{3} & x = 1 \\[6pt]
\frac{1}{x^2-1}\left(1 - \frac{1}{\sqrt{x^2-1}}\arctan\frac{\sqrt{x^2-1}}{1}\right)
& x > 1
\end{cases}
\end{equation}

In PyAutoLens, $\kappa_s = 2\rho_s r_s / \SigmaCr$ is the characteristic
convergence, so $\kappa_{\rm NFW}(x) = \kappa_s \cdot f(x)$.

% ============================================================
\section{The Bulge-Halo Conspiracy}
\label{sec:conspiracy}
% ============================================================

Observations show that the total density profile of massive ellipticals
is remarkably isothermal: $\rho_{\rm tot} \propto r^{-\gamma'}$ with
$\gamma' \approx 2.0 \pm 0.1$ (Koopmans et al.\ 2006, 2009).

This arises from a \textbf{conspiracy} between the stellar and dark
components:
\begin{itemize}[nosep]
  \item Stellar: $\rho_*(r) \propto r^{-3}$ at small $r$ (steeper than isothermal)
  \item Dark matter: $\rho_{\rm DM}(r) \propto r^{-1}$ (shallower than isothermal)
  \item Sum: the steeper-declining stars compensate the flatter dark matter,
    producing an approximately isothermal total profile
\end{itemize}

This conspiracy is not fully understood theoretically --- it may arise from
adiabatic contraction of the halo by baryons during galaxy formation.

% ============================================================
\section{Dark Matter Fraction}
\label{sec:dm-fraction}
% ============================================================

The dark matter fraction within the Einstein radius:
\begin{keyresult}
\begin{equation}
f_{\rm DM}(\thetaE) = \frac{M_{\rm DM}(<\thetaE)}{M_{\rm total}(<\thetaE)}
= 1 - \frac{M_*(<\thetaE)}{M_{\rm total}(<\thetaE)}
\label{eq:fdm}
\end{equation}
\end{keyresult}

For SLACS lenses (Auger et al.\ 2010): $f_{\rm DM}(\thetaE) \sim 0.2$--$0.5$,
increasing with Einstein radius (more massive galaxies have higher dark
matter fractions).

\subsection{The $\Upsilon_*$--$\kappa_s$ Degeneracy}

Lensing constrains the \textbf{total} mass $\kappa_* + \kappa_{\rm DM}$
at the Einstein radius. Increasing $\Upsilon_*$ (more stellar mass)
must be compensated by decreasing $\kappa_s$ (less dark matter), and
vice versa. Breaking this degeneracy requires:
\begin{itemize}[nosep]
  \item Stellar kinematics (velocity dispersion $\sigma_v$ constrains
    the mass profile slope)
  \item Stellar population synthesis (independent $\Upsilon_*$ estimate)
  \item Multiple Einstein radii (compound lenses with two source planes)
\end{itemize}

% ============================================================
\section*{References}
% ============================================================

\begin{itemize}[nosep, leftmargin=*]
  \item Auger, M.W.\ et al.\ (2010), ApJ, 724, 511
  \item Bartelmann, M.\ (1996), A\&A, 313, 697
  \item Congdon, A.B.\ \& Keeton, C.R.\ (2018), \emph{Principles of Gravitational Lensing}
  \item Koopmans, L.V.E.\ et al.\ (2006), ApJ, 649, 599
  \item Navarro, J.F., Frenk, C.S.\ \& White, S.D.M.\ (1996), ApJ, 462, 563
  \item Nightingale, J.W.\ et al.\ (2019), MNRAS, 489, 2049
  \item Treu, T.\ \& Koopmans, L.V.E.\ (2004), ApJ, 611, 739
\end{itemize}

\end{document}
